\section{Managerial Process}
\label{sec:managerial}

Our team has three members, and each owns one area: Muhammad Haider Ali is the project manager and owns the architecture, the server and the database; Zain Ali Mehmood owns the student desktop client, lockdown and monitoring; and Roshaan Ali Shah owns AI evaluation, plagiarism detection and testing. We follow Scrum in a light form \cite{scrumguide2020}: each sprint starts with planning and ends with a review with our supervisor, short written check-ins replace the daily stand-up, and a blocker unsolved after 24 hours goes to a team discussion.

\textbf{Changes since the FYP-I proposal.} The twelve proposal modules were regrouped: role management became M7, eligibility became its own module (M9), and behaviour analysis became the violation classifier used in live monitoring. We added the Invigilator and SuperAdmin roles, chose PostgreSQL on Neon as the database, and run the grading inside the server instead of in a separate Python service.

\subsection{Management Objectives and Priorities}
\label{sec:priorities}

Our objective is a system that lets a teacher run a fair lab exam from start to finish. We measure progress with four goals:
\begin{itemize}[nosep]
    \item Deliver all twelve modules, each with its own tests, within the 32 weeks.
    \item Test every functional requirement FR1--FR12 (Chapter~\ref{chap:testing}).
    \item Use no paid services.
    \item Show the supervisor a working increment at the end of every sprint.
\end{itemize}

When time is short, we protect work in this order, and lower items are cut first:
\begin{enumerate}[nosep]
    \item \textbf{Foundations:} accounts and roles, exam creation and the student dashboard (M7, M8, M3).
    \item \textbf{Exam integrity:} device binding, the readiness check, eligibility and seats, and lockdown (M1, M2, M9, M4).
    \item \textbf{No lost work:} autosave, encrypted backup and the server deadline (M6).
    \item \textbf{Live supervision:} monitoring client and proctoring dashboard (M5, M10).
    \item \textbf{Evaluation support:} test-case marking, theory scoring and plagiarism checks (M11).
    \item \textbf{Reporting:} audit log, analytics and reports (M12).
\end{enumerate}

\subsection{Assumptions and Constraints}
\label{sec:constraints}

\textbf{Assumptions}
\begin{itemize}[nosep]
    \item Lab PCs run Windows 10 or 11 (we tested on Windows 11), and we can install the client with administrator rights.
    \item Each PC reaches the server over the lab network; short outages may happen.
    \item Each student uses one assigned PC, with a webcam when the exam needs one.
    \item Teachers supply the test cases and model answers used for automatic marking.
\end{itemize}

\textbf{Constraints}
\begin{itemize}[nosep]
    \item Three members and 32 weeks over two semesters.
    \item No budget: we use the Neon free plan, and the AI features run offline on the CPU (text-similarity scoring, YuNet face detection \cite{yunet} and YOLOX-Nano phone detection \cite{yolox}). Marking with an online language model is off by default.
    \item The client runs only on Windows and cannot block Ctrl+Alt+Del.
    \item Isolated running of C and C++ answers needs Docker on the server. Without it, a local compiler that is not isolated is used, and with no compiler the answer is left for the teacher.
    \item For privacy, no face recognition is done. Screenshots and webcam pictures can be viewed only by the exam's owner, the staff assigned to supervise it and administrators, and are deleted after 30 days.
\end{itemize}
